\def\SigleNumero{STT-2902}
\def\NomCours{Modélisation statistique}
\def\DateEvaluation{10 novembre 2025}
\def\HeureEvaluation{8 h 30 à 10 h 20}
\def\Session{}
\def\NomEvaluation{Examen 2}
\def\Ponderation{}
% Ne rien mettre entre les accolades si l'enseignant (ou l'enseignante) est aussi le coordonnateur (ou la coordonnatrice)
\def\Coordonnatrice{}
\def\Coordonnateur{}

% Ne rien mettre entre les accolades s'il y a plus d'une section
\def\Enseignant{Julien Miron}
\def\Enseignante{}

% Ne rien mettre entre les accolades s'il s'agit d'un cours à section unique
% Mettre le nom de chacun des enseignants et enseignantes entre les accolades
% Une case à cocher sera générée pour chacune des sections
\def\SectionA{}
\def\SectionB{}
\def\SectionC{}
\def\SectionD{}


\def\Directives{
\begin{enumerate}
\item Matériel autorisé:
\begin{itemize}
    \item Calculatrice scientifique {\bf autorisée par la faculté des sciences et de génie};
    \item Votre cahier de notes de cours;
    \end{itemize}
\item Assurez-vous que cette évaluation comporte \numquestions ~questions réparties sur \numpages{}~pages, pour un total de \numpoints{}~points et un maximum de 8~points bonus.\
\item \textbf{Veuillez ne pas dégrafer votre examen.}
\item Ajoutez vos initiales dans le bas de chaque page.
\item Vous devez\textbf{ justifier toutes vos réponses} par une démarche claire, sauf pour les questions à choix de réponse.\
\item Pour les questions à choix de réponse et les vrais ou faux, vous pouvez mettre un X, un crochet ($\checkmark$) ou remplir la case.\
\item Le verso des feuilles peut être utilisé comme brouillon, \textbf{mais ne sera pas corrigé}. Attention à ne rien dégrafer.
\end{enumerate}
}



\documentclass{Examens/Evaluation}
\usepackage{multicol,graphicx}
\begin{document}

\pageblanche
%%%%%%%%%%%%%%%% QUESTION %%%%%%%%%%%%%%%%
\begin{questions}
\question[3] Soit la variable aléatoire $X$, dont les deux résultats possibles sont $X=1$ avec probabilité $0,6$ et $X=2$ avec probabilité $0,4$. Quelle est l'espérance de $X$ \textbf{?}

\begin{oneparcheckboxes}
    \choice 0,4
    \choice 0,6
    \choice 1
    \correctchoice 1,4
    \choice 1,6
\end{oneparcheckboxes}

\vspace{0.5cm}
\question[3] Cochez l'énoncé qui est vrai.

\begin{checkboxes}
    \choice L'espérance d'une variable aléatoire est invariante par translation, mais sa variance n'est pas invariante par translation.
    \choice L'espérance et la variance sont invariantes par translation.
    \correctchoice L'espérance d'une variable aléatoire n'est pas invariante par translation, mais sa variance est invariante par translation.
    \choice L'espérance et la variance ne sont pas invariantes par translation.
\end{checkboxes}

\vspace{0.5cm}

\question[3] Soit $X$ et $Y$, deux variables aléatoires indépendantes de variances respectives $\sigma^2_X$ et $\sigma^2_Y$. Quelle est la variance de $4X-2Y+2$ \textbf{?}

\begin{checkboxes}
    \choice $4\sigma^2_X-2\sigma^2_Y$
    \choice $4\sigma^2_X-2\sigma^2_Y+2$
    \choice $16\sigma^2_X-4\sigma^2_Y$
    \choice $16\sigma^2_X-4\sigma^2_Y+2$
    \choice $16\sigma^2_X-4\sigma^2_Y+4$
    \choice $4\sigma^2_X+2\sigma^2_Y$
    \choice $4\sigma^2_X+2\sigma^2_Y+2$
    \correctchoice $16\sigma^2_X+4\sigma^2_Y$
    \choice $16\sigma^2_X+4\sigma^2_Y+2$
    \choice $16\sigma^2_X+4\sigma^2_Y+4$
\end{checkboxes}

\vspace{0.5cm}

\question[3] Parmi les expressions suivantes, laquelle est une expression de la variance pour une variable aléatoire $X$ discète\textbf{?}

      \begin{checkboxes}
          \choice $V(X)=\int_{-\infty}^{\infty}(x-\mu)^2dx$
          \correctchoice $V(X)=E[(X-\mu)^2]$
          \choice $V(X)=\sum_{\mathbb{R}_X}(x-\mu)^2$
          \choice $V(X)=E(X)-E(X)^2$
      \end{checkboxes}
          
\vspace{0.5cm}

\question[3] La loi de Student de degrés de liberté $n-1$ a une d'espérance égale à $0$.

\begin{oneparcheckboxes}
    \correctchoice Vrai 
    \choice Faux
\end{oneparcheckboxes}

\vspace{0.5cm}

\question[3] Quel énoncé parmi les suivant est vrai\textbf{?}

\begin{checkboxes}
    \correctchoice L'intervalle de confiance pour la variance n'est pas symétrique autour de $S^2$.
    \choice L'intervalle de confiance pour l'espérance $\mu$ est symétrique autour de $\mu$.
    \choice Plus un intervalle de confiance est large, plus il est précis.
    \choice La quantité $1-\alpha$ se nomme le seuil de l'intervalle de confiance.
\end{checkboxes}

\pageblanche
\question[3] Cochez l'énoncé qui est vrai.

\begin{checkboxes}
    \choice La probabilité que $\overline{X}$ se trouve dans l'intervalle de confiance est $1-\alpha$.
    \choice $L$ et $U$ sont des variables aléatoires car ils sont impossibles à calculer.
    \correctchoice Toutes choses étant égales par ailleurs, un intervalle de confiance construit avec la variance connue (et $\sigma^2=s^2$) est plus court.
    \choice Si $[L;U]=[12,34;17,10]$ est un intervalle de confiance pour $\mu$, il nous manque d'information pour connaître $\overline{x}$.
    \choice Toutes choses étant égales par ailleurs, un intervalle de confiance construit avec $m>>n$ aura un niveau de confiance plus élevé.
\end{checkboxes}

\vspace{0.5cm}

\question[3] Cochez l'énoncé qui décrit l'erreur de type II.

\begin{checkboxes}
    \choice Il s'agit de la probabilité de ne pas rejeter $H_0$ alors que $H_0$ est fausse.
    \choice Il s'agit de la probabilité de  rejeter $H_0$ alors que $H_0$ est vraie.
    \choice Il s'agit de la décision de  rejeter $H_0$ alors que $H_0$ est vraie.
    \correctchoice Il s'agit de la décision de ne pas rejeter $H_0$ alors que $H_0$ est fausse.
\end{checkboxes}

\vspace{0.5cm} 

\question[3] Vous faites un test sur 20 observations  pour confronter les hypothèses $H_0: \mu=5$ et $H_1: \mu<5$. Vous obtenez un seuil observé (valeur P) de 0,09. Cochez l'énoncé qui est faux.

\begin{checkboxes}
\choice La moyenne échantillonnale est inférieure à 5.
\correctchoice Vous rejetez $H_0$ au seuil de 5\%.
\choice Vous rejetez $H_0$ au seuil de 10\%.
\choice Vous ne rejetez pas $H_0: \mu=4$ contre $H_1: \mu<4$ au seuil de 5\%.
\choice Vous refaites l'expérience avec 40 observations. Vous obtenez la même moyenne échantillonnale. Votre seuil observé est alors plus petit que 0,09 pour le test $H_0: \mu=5$ contre $H_1: \mu<5$.
\end{checkboxes}

\vspace{0.5cm}

\question[3] Cochez l'énoncé qui est vrai.

\begin{checkboxes}
    \choice Si $\hat{\theta}$ est un estimateur sans biais pour $\theta$, nous savons que la valeur de $\hat{\theta}$ calculé à partir d'un échantillon sera toujours très proche de $\theta$.
    \choice Toutes choses étant égales par ailleurs, pour passer d'un intervalle de confiance de niveau $95\%$ à un intervalle de confiance de niveau $99\%$, il suffit d'augmenter la taille d'échantillon.
    \choice L'erreur quadratique moyenne d'un estimateur sans biais est égale à sa variance au carré.
    \correctchoice Dans un test d'hypothèses au seuil $\alpha$, quand $H_0$ est vraie, la probabilité de ne pas faire d'erreur est égale à $1-\alpha$.
\end{checkboxes}

\pageblanche

\question[10]
Sachant que $X_1,\hdots,X_n\sim\mathcal{N}(\mu,\sigma^2)$ et $\displaystyle\frac{(n-1)S^2}{\sigma^2}\sim \chi^2_{n-1}$, où $S^2=(n-1)^{-1}\sum_{i=1}^n (X_i-\overline{X})^2$, montrez que la forme de l'intervalle de confiance de niveau $1-\alpha$ pour $\sigma^2$ est $$\displaystyle\left[\frac{(n-1)S^2}{\chi^2_{\alpha/2;n-1}};\frac{(n-1)S^2}{\chi^2_{1-\alpha/2;n-1}}\right].$$

    \pageblanche

\question Soit l'échantillon qui provient d'une variable aléatoire normale d'espérance $\mu$ et de variance $\sigma^2=0,81$:

$$2,2,2,3,3,3,4,4,4,$$
pour lequel $\displaystyle\sum_{i=1}^n x_i=27$ et $\displaystyle\sum_{i=1}^9x_i^2=87$.


\begin{parts}
    \part On souhaite faire un test d'hypothèses unilatéral à droite de seuil $2,5\%$ avec $H_0:\mu=2,7$.

    \begin{subparts}
        \subpart[2] Quelle est l'hypothèse alternative?
        
        \vspace{3cm} 

        \subpart[3] Quelle est la valeur critique du test?
    \end{subparts}
\vspace{4cm}
    \part[5] Construisez un intervalle de confiance de niveau $95\%$ pour $\mu$.
\end{parts}
\pageblanche

\Large Les quatre exercices suivants sont des exercices BONUS. Vous aurez les points d'un maximum de DEUX exercices. Pour chaque exercice, vous pouvez avoir 4 points ou 0 points. 

\normalsize

\vspace{2cm}
\bonusquestion[4] Soit $X_1,\hdots,X_m$ qui sont des variables aléatoires indépendantes et identiquement distribuées de loi Inventée de paramètre $\kappa$, qu'on note $X_i\sim Inv(\kappa)$. 
Avec cette loi, une variable aléatoire prend la valeur $2$ avec probabilité $\kappa$ et la valeur $0$ avec la probabilité $1-\kappa$.
Nous savons que $E(X_i)=2\kappa$ et $V(X_i)=4\kappa(1-\kappa)$.
Pour estimer $\kappa$, nous utilisons $$\displaystyle \hat{\kappa}=\frac{\overline{X}}{2}.$$

Quelle est la variance de $\hat{\kappa}$?


\pageblanche
\bonusquestion[4] Quelle est la longueur maximale de la marge d'erreur d'un intervalle de confiance à $95\%$ pour $p$, lorsque $n=100$? \textit{Indice: commencez par trouver pour quelle valeur de $\hat{p}$ la marge d'erreur est maximale.}

\pageblanche
\bonusquestion[4] 

Soit $X\sim Bin(n;\theta)$. Considérons le test dont les hypothèses sont $H_0:\theta =0,5$ et $H_1:\theta>0,5$. 
La règle de décision est de rejeter $H_0$ si $n$ succès sont observés (c'est-à-dire aucun échec).
Trouvez la valeur de $n$ à partir de laquelle $\alpha\leq 0,05$.

\pageblanche
\bonusquestion[4] Soit deux échantillons aléatoires indépendants $X_1,\hdots,X_n$ et $Y_1,\hdots,Y_m$. 
    Nous savons que ces échantillons proviennent respectivement de $X\sim\mathcal{N}(\mu_X,\sigma^2_X)$ et $Y\sim\mathcal{N}(\mu_Y, \sigma^2_Y)$. 
    Nous notons $S_X^2$ et $S_Y^2$ les variances échantillonnales usuelles.
    Il est possible de montrer que 
    \begin{align*}
        \frac{S_X^2/\sigma^2_X}{S_Y^2/\sigma^2_Y}\sim F_{n-1;m-1},
    \end{align*}
    où $F_{n-1;m-1}$ est la loi de Fisher avec $n-1$ degrés de liberté au numérateur et $m-1$ degrés de liberté au dénominateur.
    En utilisant la notation standard, $F_{\alpha/2;n-1;m-1}$ est le quantile tel que $P(F_{n-1;m-1}>F_{\alpha/2;n-1;m-1})=\alpha/2$.
    
    Donnez un intervalle de confiance de niveau $1-\alpha$ pour le ratio des variances $\displaystyle\frac{\sigma^2_Y}{\sigma^2_X}$.
    


\pageblanche


\begin{comment}
    

\bonusquestion[4] Supposons que les variables aléatoires $X_1,...,X_{20}$ suivent une loi exponentielle décalée de paramètre $\theta \geq0 $, dont la \textbf{fonction de densité} est donnée par 

 \begin{equation*}
  f_X(x) = \begin{cases}
         e^{-(x-\theta)},\ \text{  si  } \theta \leq x , \\
        0, \text{ sinon.}
        \end{cases}
 \end{equation*}

Nous aimerions faire un test statistique pour savoir si $\theta=0$. Nous avons donc les hypothèses

\begin{table}[ht]
    \centering
    \begin{tabular}{cc}
         $H_0: \theta=0$ & $H_1: \theta > 0$.
    \end{tabular}
\end{table}
Pour ce test, nous utiliserons la statistique
\begin{align*}
    M=\min (X_1,...,X_{20}),
\end{align*}
qui correspond à la valeur minimale de notre échantillon (nous pouvons aussi écrire $M=X_{(1)}$, dans la notation introduite au cours).
Nous savons que sous $H_0$, $M$ suit une loi exponentielle de paramètre $n=20$. Sa \textbf{fonction de répartition} est donnée par 

 \begin{equation*}
  F_M(m) = \begin{cases}
         1-e^{-20m},\ \text{  si  } 0\leq m \\
        0, \text{ sinon.}
        \end{cases}
 \end{equation*}

La table suivante donne quelques valeurs de cette fonction de répartition.
\begin{table}[ht]
    \centering
\begin{tabular}{|cc|}
  \hline
 $m$ & $F_M(m)$ \\ 
  \hline
    0,0005 & 0,01 \\
    0,0025 & 0,05 \\
    0,0053 & 0,10 \\
    0,1151 & 0,90 \\
    0,1498 & 0,95 \\ 
    0,2303 & 0,99 \\ 
  \hline
\end{tabular}
\end{table}

Nous avons l'échantillon $x_1,\hdots,x_{20}$ suivant:

\begin{table}[ht]
\centering
\begin{tabular}{rrrrrrrrrr}
  \hline
0,1613 & 0,1933 & 0,2980 & 0,3331 & 0,3423 & 0,4392 & 0,4472 & 0,4580 & 0,5289 & 0,5530 \\ 
0.5570 & 0.7737 & 0.8574 & 0.9376 & 1.2509 & 1.4475 & 1.9064 & 2.6019 & 2.8685 & 2.9829 \\ 
   \hline
\end{tabular}
\end{table}
\begin{parts}
\bonuspart Au seuil $\alpha=5\%$, quelle est la \textit{\textbf{décision}} du test. \textbf{\textit{Une décision sans justification n'obtiendra aucun point, même si la décision est bonne.}}
\vspace{3cm}
\bonuspart \textbf{Calculez} la valeur-p du test en ne gardant que 4 chiffres après la virgule. \textit{\textbf{Expliquez votre réponse!}}
\end{parts}
\end{comment}
\end{questions}

\end{document} 
